% This code is made availiable under the Creative Commons Attribution Share Alike
% License (cc-by-sa). You can find the full License text under
% [http://creativecommons.org/licenses/by-sa/3.0/]
%
% Copyright 2005 Dirk Jagdmann <doj@cubic.org>
% Copyright 2006-20013 Lutz Behnke <lutz.behnke@gmx.de>
%
\documentclass[%
   pdftex,%
   a4paper,%                                            % DIN A4
   oneside,%                                            % einseitiger Druck
   headsepline,%                                        % Strich unter der Kopfzeile
   12pt,%                                               % 12pt Schriftgr��e
]{report}
%---------------------------------------------------------------------------------------------------
%\documentclass[%
%  BCOR5mm,%                                            % Bindungskorrektur
%  pagesize,%                                           % flexible Auswahl des Papierformats
%  halfparskip,%                                        % Europ�ischer Satz: Abstand zwischen Abs�tzen
%  bibtotoc,%                                           % Literaturverzeichnis im Inhaltsverzeichnis
%]{scrreprt}                                            % KOMA-Scriptklasse Report
%---------------------------------------------------------------------------------------------------

\usepackage{hawhhcs}

\hawtyp{Masterarbeit} % default: 'Laborbericht'

\hawauthor{Lutz Behnke}

% Titel f�r Deckblatt
\hawtitel{Kooperativer Speicher: Schw�chen und Gegenma�nahmen}
%\hawtitle{Cooperative Storage: Weaknesses and Countermeasures}

% Betreuer
\hawprof{Prof.\ Dr.\ Kai v.\ Luck}
\hawkoprof{Prof.\ Dr.-Ing.\ Martin H�bner}

% Kooperation mit externer Stelle
\hawkooperation{catWorkX GmbH}

% Ort und Datum
\hawdatum{\today} % default: \today
\hawort{Hamburg} % default: 'Hamburg'

% Titel f�r Zusammenfassungsseiten (Abstract) (Deutsch und Englisch)
\hawzftitel{Entwicklung eines E-Payment-Systems auf Basis mobiler Endger�te}
\hawzftitle{Development of an ePayment-System based on mobile devices}

\hawstichworte{Netzwerk Persistenz Storage}
\hawkeywords{Network Persistence Storage}

% Zusammenfassungsseite - Zusammenfassung(de)
\hawzusammenfassung{In mehreren Bereichen der
  IT-Technik sind heute Speichersysteme nicht mehr geeignet die
  Anforderungen an Flexibilit�t, Skalierbarkeit und Stabilit�t gegen
  Fehler zu erf�llen. Kooperative Speicher scheinen ein Ausweg aus der
  Konfigurations-Falle die sich aus immer komplexeren L�sungen
  ergibt. Im folgenden wird ein Ansatz, der zuk�nftig der Gegenstand der
  Masterarbeit des Authors werden soll, in Kontext gestellt und notwendige
  Arbeiten und Untersuchungen diskutiert.}

% Zusammenfassungsseite - Zusammenfassung(Abstract)(en)
\hawabstract{(none)}%Place the english abstract here

% Titeldruck: farbig mit Logo / ohne Logo
\hawlogo{true}  %% true = farbig; false = ausdruck (ohne balken und logo)

\begin{document}
\pagenumbering{roman}
\maketitle

{\parskip=0mm \tableofcontents}
\pagenumbering{arabic}

\chapter{In the beginning}

\section{Einleitung}
\label{sec:intro}

\subsection{Das Problem}
\label{ssec:intro:problem}

Aus einer Reihe von Bereichen der IT kommen heute Forderungen nach
verbesserten Storage-Systemen. Vor allem Konzepte wie \emph{Ubiqutous
Computing} mit seiner Vielzahl von allgegenw�rtigen intelligenten
Ger�ten, \emph{Grid Computing} mit dom�nen�bergreifenden
Anwendungen und \emph{Cluster Computing} mit einer enormen Anzahl an
Rechner-Knoten verlangen nach Verbesserungen bei

\begin{itemize}
\item Ad-Hoc Konfiguration/Discovery
\item Skalierbarkeit
\item (scheinbarer) Stabilit�t
\item Organisatorische Transparenz
\end{itemize}


\chapter{More Data}

Eine gute wissenschaftliche Arbeit braucht immer \emph{mehr} Daten. Deswegen
hat diese Arbeit auch ein Kapitel mit ''more data´

�


Ein Bild sagt mehr als tausend Worte:

\includegraphics[width=0.5\textwidth]{wrong.png}

\citep{Bell76}

\bibliographystyle{dinat}
\bibliography{bibliography/references}

\hawversicherung

\end{document}

%%% Local Variables: 
%%% mode: latex
%%% TeX-master: t
%%% End: 



